\section{Buddy 分配器：\texttt{pmm\_buddy.c}}

\subsection{核心思想}

Buddy system 将物理内存分成 $2^k$ 大小的块（order $k = 0, 1, \ldots, 10$）。
每个 order 维护一个\term{空闲链表}：

\begin{figure}[H]
  \centering
  % figures/ch07/fig05-buddy-allocator-free_list-orde.tex — auto-extracted from chapter source
\begin{tikzpicture}[node distance=0cm]
  \foreach \k/\sz in {0/1, 1/2, 2/4, 3/8, 4/16} {
    \node[font=\small\ttfamily, anchor=east] (lbl\k) at (-0.5, -\k*0.7) {free\_list[\k]};
    \node[memcell, minimum width=1.5cm, fill=green!15, right=0.2cm of lbl\k] (n\k) {$2^{\k}$ 页};
    \node[memcell, minimum width=1.5cm, fill=green!15, right=0.3cm of n\k] (m\k) {$2^{\k}$ 页};
    \draw[pointer] (n\k.east) -- (m\k.west);
    \node[right=0.3cm of m\k, font=\scriptsize\ttfamily] {→ NULL};
  }
  \node[font=\small, above=0.3cm of n0] {order 0: 单页 (4K)};
  \node[font=\small, above=0.3cm of n4] {order 4: 16 页 (64K)};
\end{tikzpicture}

  \caption{Buddy allocator 的 free\_list 数组（每个 order 一个链表）}
\end{figure}

\subsection{分配：分裂（Split）}

如果请求 order $k$ 但 \texttt{free\_list[$k$]} 为空，就从更高 order 取一个块并\textbf{分裂}：

\begin{figure}[H]
  \centering
  % figures/ch07/fig06-buddy-order-2-order-1.tex — auto-extracted from chapter source
\begin{tikzpicture}[node distance=0cm]
  % order 2 block (4 pages)
  \node[memcell, minimum width=6cm, fill=cyan!20, minimum height=0.8cm] (big) {order 2 (4 页)};
  \node[above=0.2cm of big, font=\small] {从 free\_list[2] 取出};

  % 分裂成两个 order 1
  \node[memcell, minimum width=2.8cm, fill=orange!20, minimum height=0.8cm, below left=1cm and -1.5cm of big] (left) {order 1 (2 页)};
  \node[memcell, minimum width=2.8cm, fill=green!20, minimum height=0.8cm, below right=1cm and -1.5cm of big] (right) {order 1 (2 页)};

  \draw[-{Stealth[length=4pt]}, thick] (big.south west) -- (left.north) node[midway, left, font=\scriptsize] {返回给调用者};
  \draw[-{Stealth[length=4pt]}, thick] (big.south east) -- (right.north) node[midway, right, font=\scriptsize] {放回 free\_list[1]};
\end{tikzpicture}

  \caption{Buddy 分裂：order 2 块拆成两个 order 1 块}
\end{figure}

\subsection{释放：合并（Coalesce）}

释放一个 order $k$ 块后，检查它的\term{buddy}是否也空闲。如果是，合并为 order $k+1$：

\begin{lstlisting}[style=cstyle]
// 给定页帧号 P 和 order k，buddy 的帧号：
uint32_t buddy = page_frame ^ (1u << k);
\end{lstlisting}

\begin{thinkbox}
为什么 XOR 能找到 buddy？假设 order 1（2 页），页帧号 4 和 5 是一对，6 和 7 是一对。
$4 \oplus 2 = 6$（不是 buddy！）。实际上 buddy 是以\textbf{对齐的块}为单位的：
order $k$ 块的起始帧号必须是 $2^k$ 的整数倍，所以 $P \oplus 2^k$ 恰好切换到同级的另一半。
例如帧 4 的 order-1 buddy = $4 \oplus 2 = 6$，帧 6 的 buddy = $6 \oplus 2 = 4$。
\end{thinkbox}

% ============================================================================

